\begin{ex}
Tính: $35 + 47$.
\shortans{82}
\loigiai{
Ta có $35 + 47 = 82$.
}
\end{ex}

\begin{ex}
Tính: $82 - 35$.
\shortans{47}
\loigiai{
Ta có $82 - 35 = 47$.
}
\end{ex}

\begin{ex}
Tính: $48 + 36$.
\shortans{84}
\loigiai{
Ta có $48 + 36 = 84$.
}
\end{ex}

\begin{ex}
Tính: $84 - 48$.
\shortans{36}
\loigiai{
Ta có $84 - 48 = 36$.
}
\end{ex}

\begin{ex}
Tính: $59 + 25$.
\shortans{84}
\loigiai{
Ta có $59 + 25 = 84$.
}
\end{ex}

\begin{ex}
Tính: $73 - 38$.
\shortans{35}
\loigiai{
Ta có $73 - 38 = 35$.
}
\end{ex}

\begin{ex}
Tính: $60 - 24$.
\shortans{36}
\loigiai{
Ta có $60 - 24 = 36$.
}
\end{ex}

\begin{ex}
Tính: $100 - 37$.
\shortans{63}
\loigiai{
Ta có $100 - 37 = 63$.
}
\end{ex}

\begin{ex}
Tính: $100 - 65$.
\shortans{35}
\loigiai{
Ta có $100 - 65 = 35$.
}
\end{ex}

\begin{ex}
Tính giá trị biểu thức: $38 + 27 - 45$.
\shortans{20}
\loigiai{
Ta có $38 + 27 = 65$, rồi $65 - 45 = 20$.
}
\end{ex}

\begin{ex}
Tính giá trị biểu thức: $91 - 46 + 25$.
\shortans{70}
\loigiai{
Ta có $91 - 46 = 45$, rồi $45 + 25 = 70$.
}
\end{ex}

\begin{ex}
Tính giá trị biểu thức: $54 + 38 - 62$.
\shortans{30}
\loigiai{
Ta có $54 + 38 = 92$, rồi $92 - 62 = 30$.
}
\end{ex}

\begin{ex}
Điền số thích hợp: $46 + \dots = 81$.
\shortans{35}
\loigiai{
Số cần điền: $81 - 46 = 35$.
}
\end{ex}

\begin{ex}
Điền số thích hợp: $72 - \dots = 38$.
\shortans{34}
\loigiai{
Số cần điền: $72 - 38 = 34$.
}
\end{ex}

\begin{ex}
Điền số thích hợp: $\dots + 29 = 65$.
\shortans{36}
\loigiai{
Số cần điền: $65 - 29 = 36$.
}
\end{ex}

\begin{ex}
Điền số thích hợp: $\dots - 37 = 48$.
\shortans{85}
\loigiai{
Số cần điền: $48 + 37 = 85$.
}
\end{ex}

\begin{ex}
Lớp 2A có 18 bạn nam và 19 bạn nữ. Hỏi lớp 2A có bao nhiêu bạn học sinh?
\shortans{37}
\loigiai{
Cả lớp có: $18 + 19 = 37$ bạn.
}
\end{ex}

\begin{ex}
Bác nông dân nuôi 52 con gà, bác bán đi 28 con gà. Bác còn lại bao nhiêu con gà?
\shortans{24}
\loigiai{
Còn lại: $52 - 28 = 24$ con gà.
}
\end{ex}

\begin{ex}
Đoạn dây dài 65 cm, người ta cắt đi 38 cm. Đoạn dây còn lại dài bao nhiêu xăng-ti-mét?
\shortans{27}
\loigiai{
Còn lại: $65 - 38 = 27$ cm.
}
\end{ex}

\begin{ex}
Một can dầu có 40 lít, sau khi dùng còn lại 16 lít. Đã dùng bao nhiêu lít dầu?
\shortans{24}
\loigiai{
Đã dùng: $40 - 16 = 24$ lít.
}
\end{ex}

\begin{ex}
Hai số có tổng bằng 83, số thứ nhất là 45. Số thứ hai là bao nhiêu?
\shortans{38}
\loigiai{
Số thứ hai là: $83 - 45 = 38$.
}
\end{ex}

\begin{ex}
Hai số có hiệu bằng 26, số trừ là 37. Số bị trừ là bao nhiêu?
\shortans{63}
\loigiai{
Số bị trừ là: $26 + 37 = 63$.
}
\end{ex}

\begin{ex}
Hiệu của số lớn nhất có hai chữ số và số 48 là bao nhiêu?
\shortans{51}
\loigiai{
Hiệu: $99 - 48 = 51$.
}
\end{ex}

\begin{ex}
Tính tổng của số tròn chục lớn nhất có hai chữ số và số 9.
\shortans{99}
\loigiai{
Số tròn chục lớn nhất có 2 chữ số là 90. Tổng: $90 + 9 = 99$.
}
\end{ex}

\begin{ex}
Tính: $28 + 17 + 35$.
\shortans{80}
\loigiai{
Ta có $28 + 17 = 45$, rồi $45 + 35 = 80$.
}
\end{ex}
